\documentclass[11pt,a4paper]{article}
\usepackage[T1]{fontenc}
\usepackage[utf8]{inputenc}
\usepackage{lmodern}
\usepackage{geometry}
\usepackage{amsmath,amssymb}
\usepackage{listings}
\usepackage{xcolor}
\geometry{margin=25mm}

\lstdefinelanguage{Coq}{
  morekeywords={
    Inductive,Definition,Fixpoint,Lemma,Theorem,Proof,Qed,forall,exists,
    match,with,end,fun,Prop,Type,Set,apply,eapply,constructor,intros,
    inversion,subst,destruct,reflexivity,exact
  },
  sensitive=true,
  morecomment=[n]{(*}{*)}
}
\lstset{
  language=Coq,
  basicstyle=\ttfamily\small,
  keywordstyle=\color{blue!60!black},
  commentstyle=\color{green!35!black},
  breaklines=true,
  frame=single,
  columns=fullflexible
}

\title{Coq Proofs for Hindley--Milner Inference Soundness and Type Safety}
\author{}
\date{}

\begin{document}
\maketitle

\section{Goal}

This report summarizes two mechanized Coq developments about a small
Hindley--Milner-style language.

\begin{enumerate}
  \item Inference soundness: if the HM inference judgement returns a type, then
  the declarative typing judgement also assigns that type.
  \item Type safety: a closed well-typed core term is either a value or can take
  a step, and stepping preserves its type.
\end{enumerate}

The first part is implemented in \texttt{HMSoundness.v}.  The second part is
implemented in \texttt{HMTypeSafety.v}.

\section{Inference Soundness}

\subsection{Formalization}

\texttt{HMSoundness.v} defines a small expression language with variables,
integer literals, lambda abstraction, application, and \texttt{let}.  Types are
integer types, type variables, and arrow types.

\begin{lstlisting}
Inductive ty : Type :=
| TInt : ty
| TVar : nat -> ty
| TArrow : ty -> ty -> ty.
\end{lstlisting}

The declarative typing relation is called \texttt{has\_type}.  The algorithmic
HM inference relation is called \texttt{hm\_infer}.  Type schemes,
instantiation, and generalization are modeled relationally.  Thus this file
proves soundness of the inference rules with respect to declarative typing; it
does not yet prove correctness of a concrete unification algorithm.

\subsection{Main Theorem}

The central theorem is:

\begin{lstlisting}
Theorem hm_infer_sound :
  forall G e T,
    hm_infer G e T ->
    has_type G e T.
\end{lstlisting}

In mathematical notation:

\[
  \Gamma \vdash_{\mathrm{infer}} e : T
  \quad\Longrightarrow\quad
  \Gamma \vdash e : T.
\]

That is, every type produced by the inference judgement is accepted by the
declarative type system.

\subsection{Checked Examples}

The file proves that the following program has type \texttt{int}:

\[
  \texttt{let id = fun x -> x in id 42}.
\]

\begin{lstlisting}
Theorem id_program_sound :
  has_type empty id_program TInt.
\end{lstlisting}

It also proves that, assuming
\[
  + : \texttt{int -> int -> int},
\]
the expression
\[
  \texttt{fun x -> x + 1}
\]
has type \texttt{int -> int}.

\section{Type Safety}

\subsection{Operational Semantics}

\texttt{HMTypeSafety.v} defines a simply typed core language that corresponds
to the language after type inference.  It contains integers, addition, lambda
abstraction, application, and \texttt{let}.

Values are integer literals and lambda abstractions:

\begin{lstlisting}
Inductive value : tm -> Prop :=
| VInt : forall n, value (TIntLit n)
| VLam : forall x T body, value (TLam x T body).
\end{lstlisting}

The one-step relation \texttt{step} is call-by-value.  Application and
\texttt{let} reduce by substituting a closed value for a variable.

\begin{lstlisting}
Inductive step : tm -> tm -> Prop :=
| ST_AddInts :
    forall n m,
      step (TAdd (TIntLit n) (TIntLit m)) (TIntLit (n + m))
| ST_AppLam :
    forall x T body v,
      value v ->
      step (TApp (TLam x T body) v) (subst x v body)
| ST_LetValue :
    forall x v body,
      value v ->
      step (TLet x v body) (subst x v body)
| ...
\end{lstlisting}

\subsection{Substitution}

Preservation relies on the usual substitution lemma:

\begin{lstlisting}
Lemma substitution_preserves_typing :
  forall G x U t v T,
    has_type (extend G x U) t T ->
    has_type empty v U ->
    has_type G (subst x v t) T.
\end{lstlisting}

It states that replacing a variable of type \texttt{U} by a closed value of type
\texttt{U} preserves the type of the surrounding term.

The file also proves canonical forms lemmas: a value of type \texttt{int} is an
integer literal, and a value of function type is a lambda abstraction.

\subsection{Progress}

The progress theorem says that a closed well-typed term is not stuck:

\begin{lstlisting}
Theorem progress :
  forall t T,
    has_type empty t T ->
    value t \/ exists t', step t t'.
\end{lstlisting}

Equivalently:

\[
  \emptyset \vdash t : T
  \quad\Longrightarrow\quad
  \mathrm{value}(t)
  \;\lor\;
  \exists t'.\; t \to t'.
\]

\subsection{Preservation}

The preservation theorem says that one-step evaluation preserves types:

\begin{lstlisting}
Theorem preservation :
  forall t t' T,
    has_type empty t T ->
    step t t' ->
    has_type empty t' T.
\end{lstlisting}

Equivalently:

\[
  \emptyset \vdash t : T
  \land
  t \to t'
  \quad\Longrightarrow\quad
  \emptyset \vdash t' : T.
\]

Together, progress and preservation give the standard syntactic type safety
result: closed well-typed terms do not get stuck.

\section{Scope and Limitations}

The mechanization is deliberately split into two layers.

\begin{enumerate}
  \item \texttt{HMSoundness.v}: HM inference rules are sound with respect to
  declarative typing.
  \item \texttt{HMTypeSafety.v}: the typed core language is safe with respect to
  a call-by-value operational semantics.
\end{enumerate}

This is not a full verification of the MiniML implementation.  The following
parts remain future work:

\begin{itemize}
  \item correctness of a concrete unification algorithm;
  \item correctness of the occurs check;
  \item a principal type theorem;
  \item a direct bridge between polymorphic \texttt{let} and the runtime core;
  \item correspondence with the SML implementation.
\end{itemize}

Nevertheless, both the inference-soundness theorem and the progress /
preservation safety theorems are mechanically checked by Coq.

\section{Verification}

The proofs were checked with Coq 8.16.1:

\begin{lstlisting}
cd /home/yass/my-coq-project
coqc HMSoundness.v
coqc HMTypeSafety.v
\end{lstlisting}

The LaTeX report was built with:

\begin{lstlisting}
pdflatex -interaction=nonstopmode hm_safety_report.tex
\end{lstlisting}

\end{document}
